\chapter{Release 2 : Parcours client et services numériques}
\label{chap:release2}
\section*{Introduction}
Cette release concerne l'expérience client après l'arrivée à l'hôtel. Elle couvre l'espace chambre PWA, les services hôteliers numériques, les événements, le convertisseur de devises et la messagerie multilingue avec la réception.

%% ================================================================
%%  SPRINT 4
%% ================================================================
\section{Sprint 4 : Espace client PWA et accès chambre}
\label{sec:sprint4}

\subsection{Backlog du sprint}
\begin{table}[H]
\centering
\caption{Backlog du Sprint 4}
\small
\begin{tabularx}{\textwidth}{|c|X|p{2.2cm}|c|c|c|}
\hline
\textbf{ID} & \textbf{User story} & \textbf{Acteur} & \textbf{Priorité} & \textbf{Est.} & \textbf{État} \\ \hline
S4-US01 & Accéder à l'espace chambre via QR sécurisé & Client & Must & 3j & Terminé \\ \hline
S4-US02 & Naviguer dans l'espace chambre PWA installable & Client & Must & 3j & Terminé \\ \hline
S4-US03 & Valider le token de chambre côté backend & Système & Must & 2j & Terminé \\ \hline
\end{tabularx}
\end{table}

\subsection{Diagramme de cas d'utilisation}
\safefig{diagrams/sprint4_usecase.png}{Diagramme de cas d'utilisation du Sprint~4 --- Espace client PWA}{sprint4_usecase}

\subsection{Description textuelle des cas d'utilisation}
\begin{table}[H]
\centering
\caption{Description textuelle de « Accéder à l'espace chambre »}
\small
\begin{tabularx}{\textwidth}{|p{3.5cm}|X|}
\hline
Cas d'utilisation & Accéder à l'espace chambre \\ \hline
Acteur principal & Client \\ \hline
Objectif & Permettre au client d'ouvrir l'espace chambre après scan du QR. \\ \hline
Précondition & La réservation est en statut \texttt{CHECKED\_IN} et le token est valide. \\ \hline
Scénario nominal &
\begin{enumerate}[nosep]
\item Le client scanne le QR code remis par la réception.
\item Le navigateur ouvre l'URL contenant le token JWT.
\item \texttt{RoomLayout.tsx} extrait le token et appelle le backend pour validation.
\item Le backend vérifie la signature, l'expiration et le statut de la réservation.
\item L'espace chambre s'affiche avec la barre de navigation \texttt{RoomNavBar.tsx}.
\end{enumerate} \\ \hline
Scénarios alternatifs & Token expiré ou invalide : accès refusé avec message explicite. Réservation non \texttt{CHECKED\_IN} : accès refusé. \\ \hline
Postcondition & Le client accède à son espace chambre et peut naviguer vers les services. \\ \hline
\end{tabularx}
\end{table}

\subsection{Modèle conceptuel de données}
Ce sprint réutilise les entités \texttt{Room} et \texttt{Reservation} existantes. Aucune nouvelle entité n'est introduite ; la logique repose sur la validation du token JWT côté backend.

\safefig{diagrams/sprint4_mcd.png}{Modèle conceptuel de données du Sprint~4}{sprint4_mcd}

\subsection{Diagramme de séquence}
\safefig{diagrams/sprint4_sequence.png}{Diagramme de séquence du Sprint~4 --- Validation token et accès chambre}{sprint4_sequence}

\subsection{Diagramme de classes de conception}
\safefig{diagrams/sprint4_classes.png}{Diagramme de classes de conception du Sprint~4}{sprint4_classes}

\subsection{Réalisation}
L'espace client est implémenté sous les routes \texttt{/room/*}. Le composant \texttt{RoomLayout.tsx} vérifie le paramètre \texttt{token} transmis dans l'URL et charge la session client. La page \texttt{RoomHomePage.tsx} affiche les informations liées à la chambre et propose les liens vers les services disponibles. La barre de navigation \texttt{RoomNavBar.tsx} offre un accès rapide aux pages Réclamations, Messages, Services, Événements et Devise.

La PWA est configurée via le \texttt{manifest.json} et un service worker. L'installabilité concerne uniquement l'espace client chambre ; la page de check-in et les dashboards staff restent des interfaces web classiques. Le composant \texttt{PwaInstallBanner.tsx} affiche une bannière d'installation sur les appareils compatibles.

\screenshot{screenshots/room_home.png}{Page d'accueil de l'espace chambre client}{screenshot_room_home}

\subsection{Tests et validation}
Les validations portent sur les cas de token valide, token expiré et réservation non active.

\subsection{Revue et rétrospective}
Le sprint livre l'espace client chambre et son accès sécurisé. La difficulté principale concerne la gestion de l'expiration des tokens en fonction de la date de checkout.


%% ================================================================
%%  SPRINT 5
%% ================================================================
\section{Sprint 5 : Services hôteliers, événements et devise}
\label{sec:sprint5}

\subsection{Backlog du sprint}
\begin{table}[H]
\centering
\caption{Backlog du Sprint 5}
\small
\begin{tabularx}{\textwidth}{|c|X|p{2.2cm}|c|c|c|}
\hline
\textbf{ID} & \textbf{User story} & \textbf{Acteur} & \textbf{Priorité} & \textbf{Est.} & \textbf{État} \\ \hline
S5-US01 & Consulter les informations de l'hôtel & Client & Should & 2j & Terminé \\ \hline
S5-US02 & Consulter les événements publiés avec images & Client & Should & 3j & Terminé \\ \hline
S5-US03 & Utiliser le convertisseur de devises & Client & Could & 2j & Terminé \\ \hline
S5-US04 & Gérer les contenus, événements et devises & Admin/Réception & Should & 4j & Terminé \\ \hline
\end{tabularx}
\end{table}

\subsection{Diagramme de cas d'utilisation}
\safefig{diagrams/sprint5_usecase.png}{Diagramme de cas d'utilisation du Sprint~5 --- Services hôteliers}{sprint5_usecase}

\subsection{Description textuelle des cas d'utilisation}

\begin{table}[H]
\centering
\caption{Description textuelle de « Consulter les services hôteliers »}
\small
\begin{tabularx}{\textwidth}{|p{3.5cm}|X|}
\hline
Cas d'utilisation & Consulter les services hôteliers \\ \hline
Acteur principal & Client \\ \hline
Objectif & Afficher les informations et services utiles de l'hôtel. \\ \hline
Précondition & Le client dispose d'un accès chambre valide. \\ \hline
Scénario nominal & Le client ouvre la page des informations, le frontend récupère les contenus publiés via l'API publique et les affiche par type (wifi, restaurant, spa, etc.). \\ \hline
Scénarios alternatifs & Aucune information publiée : affichage d'un état vide. \\ \hline
Postcondition & Les informations de l'hôtel sont consultées. \\ \hline
\end{tabularx}
\end{table}

\begin{table}[H]
\centering
\caption{Description textuelle de « Utiliser le convertisseur de devises »}
\small
\begin{tabularx}{\textwidth}{|p{3.5cm}|X|}
\hline
Cas d'utilisation & Utiliser le convertisseur de devises \\ \hline
Acteur principal & Client \\ \hline
Objectif & Convertir un montant étranger en dinar tunisien (TND). \\ \hline
Précondition & Les taux de change sont disponibles côté backend. \\ \hline
Scénario nominal & Le client saisit un montant et sélectionne une devise source ; le système calcule et affiche le montant équivalent en TND. \\ \hline
Scénarios alternatifs & Taux absent : message d'indisponibilité. \\ \hline
Postcondition & Le montant converti est affiché. \\ \hline
\end{tabularx}
\end{table}

\subsection{Modèle conceptuel de données}
Ce sprint introduit les entités \texttt{HotelInfo} (title, content, type), \texttt{HotelEvent} (title, description, eventDate, imageUrl, imagePath, isPublished) et \texttt{CurrencyRate} (currency, rateToTnd).

\safefig{diagrams/sprint5_mcd.png}{Modèle conceptuel de données du Sprint~5}{sprint5_mcd}

\subsection{Diagramme de séquence}
\safefig{diagrams/sprint5_sequence.png}{Diagramme de séquence du Sprint~5 --- Consultation services}{sprint5_sequence}

\subsection{Diagramme de classes de conception}
\safefig{diagrams/sprint5_classes.png}{Diagramme de classes de conception du Sprint~5}{sprint5_classes}

\subsection{Réalisation}
Ce sprint ajoute les pages \texttt{RoomHotelInfoPage.tsx}, \texttt{RoomEventsPage.tsx} et \texttt{RoomCurrencyPage.tsx}. Les contenus sont gérés via \texttt{hotel.service.ts}, \texttt{event.service.ts} et \texttt{currency.service.ts}. Les images des événements sont stockées sur Supabase Storage (démonstration) via \texttt{storage.service.ts}, avec un stockage local prévu pour le déploiement privé. Les routes publiques \texttt{hotelPublic.routes.ts} et \texttt{eventPublic.routes.ts} permettent l'accès sans authentification staff.

\screenshot{screenshots/hotel_info.png}{Page des informations hôtelières}{screenshot_hotel_info}
\screenshot{screenshots/events_page.png}{Page des événements}{screenshot_events}
\screenshot{screenshots/currency_page.png}{Convertisseur de devises}{screenshot_currency}

\subsection{Tests et validation}
Les tests manuels valident l'affichage des informations, la publication/dépublication des événements et le calcul des conversions.

\subsection{Revue et rétrospective}
Le sprint livre les services numériques client non liés aux réclamations. Une amélioration possible consiste à enrichir les contenus multimédias.


%% ================================================================
%%  SPRINT 6
%% ================================================================
\section{Sprint 6 : Messagerie client-réception et traduction}
\label{sec:sprint6}

\subsection{Backlog du sprint}
\begin{table}[H]
\centering
\caption{Backlog du Sprint 6}
\small
\begin{tabularx}{\textwidth}{|c|X|p{2.2cm}|c|c|c|}
\hline
\textbf{ID} & \textbf{User story} & \textbf{Acteur} & \textbf{Priorité} & \textbf{Est.} & \textbf{État} \\ \hline
S6-US01 & Envoyer un message à la réception & Client & Should & 3j & Terminé \\ \hline
S6-US02 & Répondre au client depuis le dashboard & Réception & Should & 3j & Terminé \\ \hline
S6-US03 & Traduire automatiquement les messages & Système IA & Must & 4j & Terminé \\ \hline
\end{tabularx}
\end{table}

\subsection{Diagramme de cas d'utilisation}
\safefig{diagrams/sprint6_usecase.png}{Diagramme de cas d'utilisation du Sprint~6 --- Messagerie et traduction}{sprint6_usecase}

\subsection{Description textuelle des cas d'utilisation}

\begin{table}[H]
\centering
\caption{Description textuelle de « Envoyer un message bilingue »}
\small
\begin{tabularx}{\textwidth}{|p{3.5cm}|X|}
\hline
Cas d'utilisation & Envoyer un message bilingue \\ \hline
Acteur principal & Client \\ \hline
Objectif & Permettre au client de communiquer avec la réception dans sa langue. \\ \hline
Précondition & Le client dispose d'un accès chambre valide. \\ \hline
Scénario nominal &
\begin{enumerate}[nosep]
\item Le client écrit un message sur \texttt{RoomMessagesPage.tsx}.
\item Le backend reçoit le message via \texttt{guestMessage.service.ts}.
\item Le service appelle le microservice IA (\texttt{/detect-language} puis \texttt{/translate}).
\item La langue est détectée ; le message est traduit vers le français pour le staff.
\item Le message est stocké dans \texttt{GuestStaffMessage} avec les deux versions.
\item La réception reçoit le message en temps réel via Socket.IO.
\end{enumerate} \\ \hline
Scénarios alternatifs & Service IA indisponible : le message original est conservé sans traduction. \\ \hline
Postcondition & Le message est enregistré et visible par la réception. \\ \hline
\end{tabularx}
\end{table}

\begin{table}[H]
\centering
\caption{Description textuelle de « Répondre au client »}
\small
\begin{tabularx}{\textwidth}{|p{3.5cm}|X|}
\hline
Cas d'utilisation & Répondre au client \\ \hline
Acteur principal & Réceptionniste \\ \hline
Objectif & Envoyer une réponse compréhensible par le client. \\ \hline
Précondition & La conversation existe et le réceptionniste est authentifié. \\ \hline
Scénario nominal & La réception écrit une réponse en français ; le backend la traduit vers la langue détectée du client via NLLB-200 ; le client reçoit la réponse traduite via Socket.IO. \\ \hline
Scénarios alternatifs & Échec traduction : la réponse en français est conservée. \\ \hline
Postcondition & Le client reçoit la réponse dans sa langue. \\ \hline
\end{tabularx}
\end{table}

\subsection{Modèle conceptuel de données}
Ce sprint introduit l'entité \texttt{GuestStaffMessage} qui stocke pour chaque message : le message original, la langue détectée, la version traduite pour le staff (\texttt{staffMessage}) et la version traduite pour le client (\texttt{clientMessage}). Le type d'expéditeur est géré par l'enum \texttt{MessageSenderType} (\texttt{CLIENT} / \texttt{STAFF}).

\safefig{diagrams/sprint6_mcd.png}{Modèle conceptuel de données du Sprint~6}{sprint6_mcd}

\subsection{Diagramme de séquence}
\safefig{diagrams/sprint6_sequence.png}{Diagramme de séquence du Sprint~6 --- Message client avec traduction}{sprint6_sequence}

\subsection{Diagramme de classes de conception}
Les éléments concernés sont : \texttt{RoomMessagesPage.tsx}, \texttt{ReceptionDashboard.tsx}, \texttt{guestMessage.routes.ts}, \texttt{staffGuestMessage.routes.ts}, \texttt{guestMessage.service.ts}, \texttt{message.socket.ts}, \texttt{ai.service.ts}, \texttt{translation\_service.py}, \texttt{language\_service.py} et le modèle \texttt{GuestStaffMessage}.

\safefig{diagrams/sprint6_classes.png}{Diagramme de classes de conception du Sprint~6}{sprint6_classes}

\subsection{Réalisation}
La messagerie est disponible côté client via \texttt{RoomMessagesPage.tsx} et côté réception via l'onglet messages de \texttt{ReceptionDashboard.tsx}. Les échanges reposent sur le modèle \texttt{GuestStaffMessage} et sur Socket.IO via \texttt{message.socket.ts} pour la livraison temps réel. Le backend utilise \texttt{guestMessage.service.ts} et \texttt{ai.service.ts} pour appeler le service FastAPI.

La traduction utilise le modèle pré-entraîné \texttt{facebook/nllb-200-distilled-600M}~\cite{nllb200}, chargé en lazy loading au premier appel. Le service \texttt{translation\_service.py} implémente un cache mémoire et un glossaire hôtelier en post-traitement pour améliorer la qualité des traductions spécifiques au domaine. La détection de langue repose sur \texttt{langdetect} avec un mapping vers les codes NLLB supportés.

\screenshot{screenshots/chat_client.png}{Interface de messagerie --- Côté client}{screenshot_chat_client}
\screenshot{screenshots/chat_reception.png}{Interface de messagerie --- Côté réception}{screenshot_chat_reception}

\subsection{Tests et validation}
Les tests incluent la validation WebSocket via \texttt{websocket.test.ts}. Les tests du service IA dans \texttt{test\_language\_detection.py} vérifient la détection de langue. Les validations manuelles portent sur l'envoi client, la réception en temps réel et la réponse traduite.

\subsection{Revue et rétrospective}
Le sprint livre une messagerie bilingue temps réel. La principale difficulté concerne le coût mémoire du modèle de traduction ($\approx$~1.2~Go en RAM). Le lazy loading et le cache atténuent ce problème.

\section*{Conclusion}
Cette release a enrichi l'expérience client avec un espace chambre PWA complet : services hôteliers, événements avec images, convertisseur de devises et messagerie bilingue avec traduction automatique via NLLB-200.
